% !PRINTED 640
\weberpage{640}
%<<CONT

and obtain

\weq{5}{(\varepsilon^{r_\nu}, \xi) = K_\nu \tau_\nu^{r_{n-2}} \tau_{\nu+1}^{r_{n-3}} \ldots \tau_{\nu+n-2}^{r_0},}

and if one adds all these formulae, and further sets the rational quantity $(1, x) = A$:

\weq{6}{n\xi_0 = A + \sum_{0,\,n-2}^{\nu} K_\nu \tau_\nu^{r_{n-2}} \tau_{\nu+1}^{r_{n-3}} \ldots \tau_{\nu+n-2}^{r_0}.}

By \S.\,184, (16) and (17) the quantities $K_\nu$ and $k_\nu$ are all different from zero.

By (6) one of the roots $\xi$ is represented.

If we set, for abbreviation,

\[
R_\nu = k_\nu^{r_{n-2}} k_{\nu+1}^{r_{n-3}} \ldots k_{\nu+n-2}
\]

and write (6) in the form

\weq{7}{n\xi_0 = A + K_0 \sqrt[n]{R_0} + K_1 \sqrt[n]{R_1} + \cdots + K_{n-2} \sqrt[n]{R_{n-2}},}

then $\xi_0$ is thus expressed by $n-1$ radicals of degree $n$, each of which can in itself have $n$ different values. But these values cannot be independent of one another, since otherwise the number of values of $\xi$ that result from (7) would be too great. One obtains in fact from (5), if by \S.\,184, (16) one sets

\[
(\varepsilon^{r_\nu}, \xi) (\varepsilon^{r_{\nu-1}}, \xi)^{-g} = k_{\nu-1}
\]

\weq{8}{K_\nu \sqrt[n]{R_\nu} = K_{\nu-1}^g k_{\nu-1} \left( \sqrt[n]{R_{\nu-1}} \right)^g,}

and can accordingly express all these radicals rationally by one of them and by $k_\nu$.

But the expression (6) has over (7) the great advantage \gesperrt{that, however one may determine the $n-1$ radicals $\sqrt[n]{k_\nu}$, it nevertheless represents only $n$ different values, namely the $n$ roots $\xi$}.

To prove this, we denote by $\varepsilon_0, \varepsilon_1 \ldots \varepsilon_{n-2}$ any arbitrary system of $n$th roots of unity, and replace in (5)

\[
\sqrt[n]{k_0}, \quad \sqrt[n]{k_1} \ldots \sqrt[n]{k_{n-2}}
\]

by

\[
\varepsilon_0 \sqrt[n]{k_0}, \; \varepsilon_1 \sqrt[n]{k_1} \ldots \varepsilon_{n-2} \sqrt[n]{k_{n-2}}:
\]
%CONT>>
